\subsection{From Individual Feedback to Local Credit}
Token-wise updates follow each teacher preference independently; block-wise updates aggregate the preferences and update the expression's components with one signal. In Figure~\ref{fig:credit}, averaging $(2,-1,2)$ gives positive feedback at all three positions. This couples adjacent updates. A fixed-width block is an operational unit, not an annotated reasoning step.

The student generates a complete response; the frozen teacher scores its prefixes; detached feedback is aggregated locally; and gradients update only the current student. No new teacher responses, reward model or process labels are needed. Teacher scoring can process the generated sequence in a batched forward pass, but grouping does not reduce the number of conditional probabilities scored.

\subsection{Block-Level On-Policy Distillation}
\paragraph{Shared feedback and joint ratio.}
Group consecutive response tensor positions into disjoint blocks $B_j$ of width $k$, without packing across mask holes. Below, sums within a block select valid positions, $n_j$ is their count, and $T=\sum_j n_j$. Empty blocks have zero weight. The ideal real-arithmetic reduction, omitting the implementation's $10^{-8}$ denominator stabilizer, is
\begin{align}
 \bar a_j&=\frac{1}{\max(1,n_j)}\sum_{t\in B_j}a_t,
 &R_j(\theta)&=\exp\left(\sum_{t\in B_j}\log r_t(\theta)\right),\label{eq:block}\\
 \mathcal L_{\mathrm{block}}&=-\frac1T\sum_{j:n_j>0}n_j\mathcal C(R_j,\sg(\bar a_j)),
 &\mathcal L_{\mathrm{token}}&=-\frac1T\sum_{t=1}^T\mathcal C(r_t,\sg(a_t)).\label{eq:loss}
\end{align}
\method{} uses $k=3$. Its joint ratio $R_j=\prod_{t\in B_j}r_t$ measures the product of policy changes across the sampled block, rather than their average. Valid-length weighting gives a two-token tail weight $2/3$ relative to a full block. The code implements the same reduction using fractional block masks $n_j/k$; Appendix~\ref{app:proofs} gives the exact stabilized denominator and micro-batch convention.

\paragraph{Optimization.}
Each collected batch supplies shared advantages and current-to-old block ratios for the clipped update. This uses sampled blocks, not enumeration of a $|V|^k$ distribution or optimizer steps after every three tokens. Relative to Token OPD, the intervention changes advantage sharing, ratio granularity and loss reduction together.

\begin{figure}[t]
\centering
\includegraphics[width=.94\linewidth]{credit_assignment_concept.png}
\caption{Illustrative feedback sharing, not an observed rollout. Averaging $(2,-1,2)$ gives $(1,1,1)$, reversing the middle token's local signal. The implemented method additionally changes ratios and loss reduction. A local sign reversal does not by itself establish an error in task-level credit assignment.}
\label{fig:credit}
\end{figure}

\subsection{Understanding the Block Update}
\paragraph{Cross-token credit.}
Consider equal-sized blocks, $T=kJ$, at the old policy, where every ratio is one and clipping is inactive. Let $g_t=\nabla_\theta\log\pi_\theta(y_t\mid h_t)$. The policy-gradient components satisfy
\begin{align}
 \nabla\mathcal L_{\mathrm{token}}&=-\frac1T\sum_t a_tg_t,\nonumber\\
 \nabla\mathcal L_{\mathrm{block}}&=-\frac1T\sum_j\sum_{i\in B_j}\sum_{t\in B_j}a_i g_t
 =\nabla\mathcal L_{\mathrm{token}}-\frac1T\sum_j\sum_{\substack{i,t\in B_j\\i\ne t}}a_i g_t.
 \label{eq:cross}
\end{align}
The additional terms send token $i$'s signal to token $t$'s score gradient. They can align or conflict with the baseline; their sign alone does not tell us whether downstream mathematical reasoning benefits.

\begin{proposition}[Equal feedback does not imply equal updates]
At $\theta=\theta_{\mathrm{old}}$, with inactive clipping, full blocks of size $k$, and $a_t=a$ for all valid tokens, both policy losses equal $-a$, but
$\nabla\mathcal L_{\mathrm{block}}=k\nabla\mathcal L_{\mathrm{token}}$.
\end{proposition}
Substituting $\bar a_j=a$ in Equation~\ref{eq:loss} proves the statement. The factor concerns the policy-loss derivative, not the subsequent Adam update after gradient clipping and other terms. It reveals a scale change even when all local feedback agrees.

\paragraph{Clipping geometry.}
Away from the old policy, the joint ratio also changes which updates are clipped. Three token ratios of 1.1 produce a block ratio of $1.1^3=1.331$: each token can remain inside a conventional 0.2 clipping band while the block leaves it. Learning-rate rescaling alone therefore cannot generally recover the token-level update.

\paragraph{When averaging could help.}
To separate the averaging hypothesis from the implementation above, consider a hypothetical token-ratio, token-normalized intervention using $Wa$ instead of $a$, where $W$ is deterministic and each row sums to one. Suppose, only as a local signal model, $a=\mu+\varepsilon$, $\E\varepsilon=0$, and $\Cov(\varepsilon)=\Sigma$. Then
\begin{equation}
 \E\|Wa-\mu\|_2^2=\|(W-I)\mu\|_2^2+\operatorname{tr}(W\Sigma W^\top).
 \label{eq:mse}
\end{equation}
Averaging helps this estimator when noise reduction exceeds the bias introduced by combining different latent signals. With a constant signal and independent equal-variance noise within a block, the per-position variance becomes $\sigma^2/k$; correlated noise or heterogeneous signals can remove this benefit. This explains the motivation for short-range sharing, but is a statement about signal estimation, not a guarantee of higher task reward. The implemented objective additionally changes ratios and reduction. Appendix~\ref{app:proofs} provides the derivation, a counterexample to unconditional smoothing benefit, and a restricted interpretation in which block cross terms recover within-block future credit.
